\section*{Project contract}
Implement \texttt{pack\_rgb565} and \texttt{unpack\_rgb565}. Each receives a full
source address in \texttt{a0}, a full destination address in \texttt{a1}, and a
pixel count from zero through 32 in \texttt{a2}. The source and destination
allocations are disjoint. Both slices may start at any byte alignment. Return the
original pixel count in \texttt{v0}, preserve every bit of \texttt{s0} through
\texttt{s7}, \texttt{fp} and \texttt{sp}, and return through the supplied
\texttt{ra} without console output. Argument and temporary registers may change;
vector registers and their overlapping floating-point register storage may change.

Packed RGB888 uses consecutive unsigned red, green and blue bytes, three per pixel.
RGB565 uses a sixteen-bit value with red in bits 15 through 11, green in bits 10
through 5 and blue in bits 4 through 0. Store its low byte first. Packing retains
the high five red bits, high six green bits and high five blue bits, discarding
the remaining bits. Expansion uses the bit-repetition rule defined below. There
is no alpha component, row padding, color-space conversion or overlap handling.

Packing has a logical source slice of three times count bytes and a destination
slice of two times count bytes. Expansion exchanges those sizes. Repeated accesses
inside a permitted slice are allowed; no exact read or write count is required. Reads may also inspect bytes inside the selected logical destination.
No access may cross a logical slice boundary, even when the surrounding allocation
is mapped. Zero count performs no caller-data access. Only the logical destination
may change, including temporary writes. Guards, input, other buffers and caller
metadata remain unchanged. Submission-owned constants and private native stack
storage are permitted. Equivalent scalar and vector implementations are accepted
when they preserve this complete interface.
